\makeabstractSHORT
{ % #1: Title
Singular Fibers of Bending Systems on Hexagon Spaces
}
{ % #2: Authorship
Arsenii Zharkov\textsuperscript{1}
}
{ % #3: Affiliations (use \\ to start a newline)
\textsuperscript{1} Department of Mathematics, Amherst College, Amherst, MA
}
{ % #4: Abstract -- NOTE: REQUIRES AMSSYMB PACKAGE
The moduli space $\mathcal P(r_1,\dots,r_n)$ of closed $n$-gons in $\mathbb R^3$ with fixed side lengths $r_i>0$ modulo rigid rotations is a smooth $2(n-3)$-dimensional symplectic manifold for generic side lengths. Hamiltonian actions generated by diagonal bending flows yield completely integrable systems on these spaces, indexed by triangulations of the $n$-gon. For $n=6$, there are two combinatorially distinct integrable systems: the caterpillar (fan) system $F_{\text{cat}} = (l_{12}, l_{56}, l_{123})$ and the symmetric (star) system $F_{\text{sym}} = (l_{12}, l_{34}, l_{56})$. The image of the moment map $F: \mathcal{P}(r) \to \mathbb{R}^3$ is a 3D convex moment polytope.

A vertex of the moment polytope is Delzant if it is simple and smooth. At such points, the action is locally toric and the fiber is a single point. When a vertex is non-Delzant, the corresponding singular fiber can be a higher-dimensional manifold. We classify all conditions on generic side length vectors under which non-Delzant vertices appear for both bending systems and explicitly compute the topology of their singular fibers.

In the symmetric system, a non-Delzant vertex occurs if and only if two adjacent side lengths are equal (e.g., $r_1 = r_2$) and the others are at their extremes. Reconstruction reveals that every non-Delzant vertex in $F_{\text{sym}}$ has a singular fiber diffeomorphic to $S^2$. In the caterpillar system, non-Delzant vertices arise either from an edge equality as above or when $z = l_{123} = 0$, causing the hexagon to decouple into two independent closed triangles $T_1 = (r_1,r_2,r_3)$ and $T_2 = (r_4,r_5,r_6)$. We prove that $F_{\text{cat}}^{-1}(v) \cong SO(3)$ if and only if both triangles strictly satisfy the triangle inequalities in this latter case. In all other generic non-Delzant cases, including degenerate triangles, the fiber is diffeomorphic to $S^2$. This completely classifies the topology of singular fibers over non-Delzant vertices in generic hexagon spaces.
}
 \newpage